% 04-ludification.tex — État de l'art 3/5 : ludification et jeux sérieux (≈ 4 pages).

\chapter{Ludification et jeux sérieux}
\label{chap:ludification}

\section{Définitions}

\begin{arediger}{1 page}
Ludification (éléments de jeu dans un contexte non ludique,
\cite{deterding2011}) contre jeu sérieux ; mécaniques usuelles : objectifs,
paliers, rétroaction, défi adapté.
\end{arediger}

\section{Effets sur la motivation et l'apprentissage}

\begin{arediger}{2 pages}
Résultats contrastés de la littérature : motivation à court terme, risque de
\enquote{chasse au curseur} sans compréhension. Condition d'efficacité : la
règle du jeu doit être la relation physique elle-même. C'est le principe de
notre mode jeu : gagner exige \(v_0^2 \geq 5 g R\) et \(m (v_0^2/R + g) \leq T_{\max}\).
\end{arediger}

\section{Réalisme et sobriété}

\begin{arediger}{1 page}
Arbitrage entre animation spectaculaire et sobriété : la première version du
jeu (personnage animé, confettis) a été remplacée par une version sobre
(\enquote{réussi} / \enquote{raté}) mais physiquement plus réaliste (corde
souple, eau en particules). Discuter ce choix à la lumière du principe de
cohérence \parencite{mayer2009}.
\end{arediger}
